\documentclass[11pt,a4paper]{article}

\usepackage{iftex}
\ifPDFTeX
  \usepackage[utf8]{inputenc}
  \usepackage[T2A]{fontenc}
  \usepackage[russian]{babel}
\else
  \usepackage{fontspec}
  \usepackage{polyglossia}
  \setmainlanguage{russian}
  \setmainfont{DejaVu Serif}
  \setsansfont{DejaVu Sans}
  \setmonofont{DejaVu Sans Mono}
\fi

\usepackage[a4paper,margin=2.1cm]{geometry}
\usepackage{xcolor}
\usepackage{hyperref}
\usepackage{tabularx}
\usepackage{booktabs}
\usepackage{longtable}
\usepackage{array}
\usepackage{enumitem}
\usepackage{titlesec}
\usepackage{fancyhdr}
\usepackage{setspace}
\usepackage{ragged2e}
\usepackage{colortbl}

\hypersetup{
  colorlinks=true,
  urlcolor=blue!55!black,
  linkcolor=blue!55!black,
  pdftitle={AI-культура. Занятие 2. Полный пакет},
  pdfauthor={OpenAI}
}

\setstretch{1.08}
\setlength{\headheight}{14pt}
\pagestyle{fancy}
\fancyhf{}
\lhead{AI-культура: трек «AI аналитика данных»}
\rhead{Занятие 2}
\cfoot{\thepage}

\definecolor{Accent}{HTML}{274C77}
\definecolor{Soft}{HTML}{EEF4FA}
\definecolor{SoftGray}{HTML}{F3F5F7}
\definecolor{Warn}{HTML}{FFF4E5}

\titleformat{\section}{\Large\bfseries\color{Accent}}{\thesection.}{0.5em}{}
\titleformat{\subsection}{\large\bfseries\color{Accent}}{\thesubsection.}{0.5em}{}

\newcolumntype{Y}{>{\RaggedRight\arraybackslash}X}
\newcommand{\badge}[1]{\colorbox{Soft}{\strut\hspace{0.35em}\textbf{#1}\hspace{0.35em}}}

\begin{document}

\begin{center}
{\LARGE \textbf{Занятие 2. Введение в анализ данных. Где живут данные?}}\\[0.4em]
{\large Магистратура, трек «AI аналитика данных»}\\[0.6em]
{\normalsize Полный пакет материалов для лектора и студента}
\end{center}

\vspace{0.6em}

\noindent\colorbox{Soft}{\parbox{\dimexpr\linewidth-2\fboxsep}{
\textbf{Главная идея занятия.}
Это не обзор ресурсов и не раннее погружение в pandas. Пара учит студента принимать первое аналитическое решение: \textbf{как найти данные, проверить их происхождение, метаданные, лицензию и ограничения до начала анализа}.
}}

\section{Место занятия в треке}

Согласно описанию трека, занятие должно познакомить студентов с open data, источниками данных (Kaggle, поисковики датасетов, гос.\ порталы, API), структурой датасета и правовыми аспектами использования данных, а практическим результатом становится \textbf{паспорт найденного датасета}. Именно поэтому занятие выступает мостом между вводной встречей про ИИ-инструменты и последующими лабораторными по очистке, преобразованию и визуализации.

\section{Ограничения и рабочие допущения}

\begin{itemize}[leftmargin=1.2em]
\item Поток может достигать \textbf{70 студентов одновременно}.
\item Уровень владения кодом у студентов разный, поэтому занятие должно одинаково работать в режимах \textbf{code} и \textbf{no-code}.
\item Преподавательский пакет должен позволять провести пару даже лектору, который понимает логику анализа данных, но не является автором материала.
\item Основной очный фокус --- на практике; лекционная часть даётся компактно, но внятно и связно.
\end{itemize}

\section{Что должно получиться у студента}

К концу пары студент:
\begin{enumerate}[leftmargin=1.3em]
\item различает \textbf{датасет}, \textbf{портал}, \textbf{API}, \textbf{формат файла};
\item умеет найти как минимум \textbf{три} источника данных по теме;
\item может прочитать базовые метаданные: владелец, дата обновления, лицензия, формат, описание полей;
\item понимает, что такое \textbf{единица наблюдения};
\item объясняет, почему выбранный источник \textbf{пригоден} или \textbf{непригоден} для аналитической задачи;
\item оформляет \textbf{паспорт датасета}.
\end{enumerate}

\section{Состав пакета материалов}

\begin{longtable}{>{\raggedright\arraybackslash}p{0.23\linewidth} >{\raggedright\arraybackslash}p{0.22\linewidth} >{\raggedright\arraybackslash}p{0.47\linewidth}}
\toprule
\textbf{Файл} & \textbf{Для кого} & \textbf{Назначение} \\
\midrule
\endhead
\textbf{Master guide (tex/pdf)} & Лектор / команда & Основной методический документ: логика занятия, тайминг, активности, критерии, работа с потоком 70 человек.\\
\textbf{Lecturer cheatsheet} & Лектор & Короткая шпаргалка на одну встречу.\\
\textbf{Instructor notebook} & Лектор & Colab-ноутбук с slide-like блоками, демонстрациями CSV/JSON/API, подсказками и комментариями для ведения пары.\\
\textbf{Student notebook} & Студент & Рабочий Colab-лист: мини-примеры, задания, шаблоны ответов, optional code track.\\
\textbf{Student longread} & Студент / платформа & Короткий текст по ключевым понятиям, который можно выдать до или после пары.\\
\textbf{No-code workbook} & Студент & No-code шаблон: сравнение 3 источников, чек-лист, паспорт датасета.\\
\textbf{HTML passport builder} & Студент & Лёгкий автономный HTML-конструктор паспорта датасета.\\
\textbf{Program/cluster matrix} & Команда курса / лектор & Матрица направлений магистратуры, кластеры потока и стартовые открытые источники по тематикам.\\
\bottomrule
\end{longtable}

\section{Рекомендуемый сценарий пары (90--100 минут)}

\rowcolors{2}{white}{SoftGray}
\begin{tabularx}{\linewidth}{>{\raggedright\arraybackslash}p{0.12\linewidth} >{\raggedright\arraybackslash}p{0.28\linewidth} Y}
\toprule
\textbf{Время} & \textbf{Блок} & \textbf{Что происходит} \\
\midrule
0--10 мин & Вход в тему & Кейсовый вопрос: \emph{«Вам дали тему анализа, но не дали данные. Как не схватить мусор?»}. Кратко фиксируются цель пары и конечный артефакт.\\
10--25 мин & Мини-лекция & Open Data, dataset vs.\ portal vs.\ API, единица наблюдения, метаданные, лицензия, правовые ограничения, роль LLM как помощника, а не замены проверки.\\
25--40 мин & Демо преподавателя & Один CSV, один JSON, один API; затем быстрый реальный поиск источника и разбор метаданных.\\
40--75 мин & Практика в мини-группах & Группы по 2--3 человека ищут 3 источника по своей теме/кластеру, сравнивают их и оформляют паспорт 1 выбранного датасета.\\
75--90 мин & Мини-защиты & 4--6 коротких выступлений: что выбрали, почему, какой главный риск/ограничение заметили.\\
+10 мин (опц.) & Буфер / домашнее задание & Разбор типичных ошибок, объяснение домашнего хвоста и критериев сдачи.\\
\bottomrule
\end{tabularx}

\section{Лекционная часть: насколько подробно её разрабатывать}

Лекционный блок здесь нужен, но \textbf{не как длинная автономная лекция}. Оптимальная плотность --- 5 компактных смысловых узлов:

\begin{enumerate}[leftmargin=1.3em]
\item \textbf{Что такое open data и почему аналитик начинает именно с поиска источника.}
\item \textbf{Где живут данные:} каталог датасетов, тематический портал, гос.\ портал, API.
\item \textbf{Как выглядит структура данных:} таблица, JSON, документация, описание полей.
\item \textbf{Что читать в метаданных:} publisher, дата обновления, лицензия, coverage, формат.
\item \textbf{Правовые и этические риски:} персональные данные, ограничения лицензии, отсутствие описания, невозможность воспроизводимости.
\end{enumerate}

Лекция должна быть достаточно подробной, чтобы любой понимающий лектор смог по ней вести пару, но не настолько тяжёлой, чтобы она съела практику.

\section{Как вести занятие на потоке около 70 человек}

\begin{itemize}[leftmargin=1.2em]
\item Не отправлять всех в полностью свободный поиск. Лучше заранее раздать \textbf{кластеры тематики} и \textbf{стартовые источники}.
\item Работать в \textbf{парах/тройках}. Это резко снижает хаос и число вопросов.
\item Основной артефакт у всех один и тот же: \textbf{паспорт датасета}.
\item Для студентов без кода дать полноценный no-code путь: браузер + workbook + HTML builder.
\item API показывать всем как идею, но делать обязательным только как \textbf{демо} или \textbf{опциональный трек}.
\end{itemize}

\section{Code / no-code логика занятия}

\subsection{Обязательный общий слой}
\begin{itemize}[leftmargin=1.2em]
\item поиск и сравнение источников;
\item анализ метаданных;
\item проверка лицензии и ограничений;
\item заполнение паспорта датасета.
\end{itemize}

\subsection{No-code путь}
\begin{itemize}[leftmargin=1.2em]
\item браузер;
\item Google Dataset Search / Kaggle / data.gov.ru / тематические порталы;
\item Google Sheets или \textbf{No-code workbook};
\item \textbf{HTML passport builder}.
\end{itemize}

\subsection{Code/light-code путь}
\begin{itemize}[leftmargin=1.2em]
\item Colab;
\item чтение CSV;
\item просмотр структуры JSON;
\item один GET-запрос к открытому API без авторизации;
\item первичная интерпретация полей.
\end{itemize}

\section{Что считать результатом пары}

\noindent\colorbox{Warn}{\parbox{\dimexpr\linewidth-2\fboxsep}{
\textbf{Главный артефакт: паспорт датасета.}
Нужно проверять не то, ``нашёл ли студент ссылку'', а то, \textbf{понял ли он, что за данные он собирается использовать}.
}}

\vspace{0.5em}

Минимальные поля паспорта:
\begin{itemize}[leftmargin=1.2em]
\item название набора;
\item ссылка на источник;
\item тип источника (портал / каталог / API / файл);
\item publisher / владелец;
\item дата публикации или обновления;
\item лицензия / условия использования;
\item формат;
\item единица наблюдения;
\item ключевые поля;
\item краткое описание возможной аналитической задачи;
\item ограничения и риски.
\end{itemize}

\section{Домашнее продолжение}

Домашка должна \textbf{продолжать} работу с тем же набором данных, а не создавать новую задачу:
\begin{enumerate}[leftmargin=1.3em]
\item довести паспорт до финальной версии;
\item добавить краткое описание полей;
\item подтвердить, каким образом данные будут импортированы (CSV/JSON/API/download link);
\item сформулировать 2--3 возможные аналитические гипотезы;
\item optional: сделать первый импорт в Colab или предпросмотр в таблице.
\end{enumerate}

\section{Быстрая рубрика оценки}

\rowcolors{2}{white}{SoftGray}
\begin{tabularx}{\linewidth}{>{\raggedright\arraybackslash}p{0.33\linewidth} >{\centering\arraybackslash}p{0.12\linewidth} Y}
\toprule
\textbf{Критерий} & \textbf{Вес} & \textbf{Что считается хорошим результатом} \\
\midrule
Релевантность источника & 25\% & Источник действительно помогает отвечать на заявленный вопрос.\\
Чтение метаданных & 25\% & Заполнены publisher, дата, лицензия, формат, описание полей/структуры.\\
Критическое мышление & 25\% & Студент видит ограничения, риски и не переоценивает качество данных.\\
Ясность паспорта & 25\% & Артефакт читается быстро, структурно и воспроизводимо.\\
\bottomrule
\end{tabularx}


\section{Стартовые тематические кластеры для большого потока}

Ниже --- \textbf{сокращённая} версия кластеров. Полный перечень программ и источников вынесен в workbook \textit{lesson02\_program\_clusters\_and\_sources.xlsx}.

\rowcolors{2}{white}{SoftGray}
\begin{tabularx}{\linewidth}{>{\raggedright\arraybackslash}p{0.31\linewidth} >{\raggedright\arraybackslash}p{0.30\linewidth} Y}
\toprule
\textbf{Кластер} & \textbf{Примеры тем} & \textbf{Стартовые источники} \\
\midrule
AI, Data Science и разработка ПО & ML/AI, open-source, software engineering, developer ecosystem & Stack Overflow Survey; GH Archive; OpenAlex; Kaggle. \\
Кибербезопасность, инфраструктура и сервисы & уязвимости, сетевые измерения, инфраструктура & CISA KEV; NVD; RIPE Atlas; CAIDA. \\
Био, здоровье, химия и foodtech & здоровье, biotech, химия, population indicators & WHO GHO; WHO GHE; World Bank; OpenAlex. \\
Урбанистика, экология, климат и энергетика & городская среда, климат, устойчивость, энергия & data.gov.ru; NOAA CDO; Data.gov; World Bank. \\
Инженерия, робототехника, фотоника и производство & сенсоры, робототехника, измерения, производство & NASA Open Data; NOAA/NCEI; UCI ML Repository; Data.gov. \\
Бизнес, продукт, дизайн и гуманитарные исследования & экономика, продукт, медиа, культура, humanities & World Bank; ECB; Europeana; Google Trends; OpenAlex. \\
\bottomrule
\end{tabularx}

\section{Какие вопросы задавать студентам вслух}

\begin{itemize}[leftmargin=1.2em]
\item Кто собрал эти данные и зачем?
\item Что является единицей наблюдения?
\item Когда датасет обновлялся в последний раз?
\item Есть ли у него лицензия или хотя бы внятные условия использования?
\item Можно ли по этим данным ответить на ваш вопрос, или данных недостаточно?
\item Что в этом наборе вызывает недоверие?
\end{itemize}

\section{Что не стоит делать на этом занятии}

\begin{itemize}[leftmargin=1.2em]
\item делать длинный блок про программирование;
\item заставлять весь поток проходить через один и тот же сложный API-клиент;
\item проверять только «нашли 3 ссылки» без обсуждения качества;
\item превращать правовой блок в изолированную сухую лекцию;
\item перегружать студентов десятком платформ вместо повторяющегося сценария поиска и оценки.
\end{itemize}

\section{Финальная рекомендация}

Лучший формат для этого занятия --- \textbf{единый педагогический сценарий + две технические дорожки}. Все студенты делают одну и ту же аналитическую работу: \textbf{поиск, сравнение, выбор и описание источника данных}. Отличается только инструментальный слой: кто-то делает это через браузер и workbook, кто-то --- через Colab и лёгкий код.

\end{document}
